% ============================================================
% §8a  Zunifikowana akcja TGP — wyprowadzenie κ z wariacji
% ============================================================
\subsection{Zunifikowana akcja TGP i~wyprowadzenie $\kappa$}%
\label{ssec:unified-action}

\statuslabel{Propozycja}

% =====================================================================
% G.0 CLOSURE 2026-05-02 (Phase 4 implementation, v1.x -> v2.0)
% =====================================================================
% Ta sekcja istnieje w dwoch warstwach:
%
%   v1.x (lin. ponizej, az do konca podsekcji): historyczna formulacja
%        z V_orig(psi) = (beta/3)psi^3 - (gamma/4)psi^4 + sqrt(-g)=c0*psi
%        (forma I metryki M9.1, FALSIFIED 2026-04-25). Zachowana
%        jako derivation chain dla sledzenia historii i diff.
%
%   v2.0 (G.0 ADDENDUM na koncu podsekcji, nowe propositions
%        prop:V-M911-canonical, prop:psi-EOM-R3, prop:vacuum-stability-G0,
%        prop:kappa-corrected-G0): kanoniczna G.0-closed formulacja
%        z V_M911(psi) = -gamma*psi^2*(4-3*psi)^2/12 + sqrt(-g)=c0*psi/(4-3*psi)
%        (M9.1'' canonical), source coef 5q/Phi_0, q*c^2/Phi_0=(4/5)pi*G_0.
%
% Wartosc kappa = 3/(4*Phi_0) jest INVARIANT po re-fit Phi_0 (P32 sympy LOCK).
% Wszystkie observables (Newton G_0, PPN gamma=beta=1, BBN, LLR, CMB,
% mass spectrum lepton z PDG <0.01%) IDENTYCZNE w v1.x i v2.0 po re-fit.
%
% Patrz:
%   research/op-g0-r3-from-canonical-projection/README.md
%   research/op-g0-r3-from-canonical-projection/sek08a_v2_specification.md
%   research/op-g0-r3-from-canonical-projection/Phase{1,2,3}_results.md
% =====================================================================

Niniejsza podsekcja realizuje trzy cele:
\begin{enumerate}[(i)]
  \item Konsoliduje rozproszone cz\l{}ony dzia\l{}ania TGP
    (hipoteza~\ref{hyp:action}, stw.~\ref{prop:ghost-free-fundamental},
    tw.~\ref{thm:D-uniqueness}) w~\textbf{jedn\k{a} zunifikowan\k{a}
    form\k{e}} $S_{\mathrm{TGP}}$.
  \item Wyprowadza r\'ownanie pola~\eqref{eq:full-field-alt}
    z~warunku $\delta S_{\mathrm{TGP}}/\delta\Phi = 0$.
  \item Pokazuje, \.ze poprawne wyprowadzenie sta\l{}ej sprz\k{e}\.zenia
    kosmologicznego $\kappa$ z~akcji daje
    $\kappa = 3/(4\PhiZero)$, a~\textbf{nie} $3/(2\PhiZero)$~---
    rozwi\k{a}zuj\k{a}c napi\k{e}cie N0-7.
\end{enumerate}

% =============================================================
\subsubsection{Pe\l{}na akcja TGP}\label{sssec:full-action}
% =============================================================

\begin{hypothesis}[Zunifikowana akcja TGP]\label{hyp:unified-action}
Ca\l{}kowita akcja TGP ma posta\'c:
\begin{equation}\label{eq:S-TGP-unified}
  \boxed{
    S_{\mathrm{TGP}}[\Phi,\psi_m]
    = \int d^4x\;\sqrt{-g_{\mathrm{eff}}}\;\biggl[
      \mathcal{L}_{\mathrm{field}}(\Phi)
      + \mathcal{L}_{\mathrm{mat}}(\Phi,\psi_m)
    \biggr],
  }
\end{equation}
gdzie:
\begin{enumerate}[(a)]
  \item \textbf{Lagran\.zjan pola} (sektor generacji przestrzeni):
  \begin{equation}\label{eq:L-field}
    \mathcal{L}_{\mathrm{field}}
    = \frac{1}{2}\,K(\varphi)\,
      g_{\mathrm{eff}}^{\mu\nu}\,
      \partial_\mu\varphi\,\partial_\nu\varphi
    \;-\; V(\varphi),
  \end{equation}
  z~$\varphi = \Phi/\PhiZero$ (zmienna bezwymiarowa), sprz\k{e}\.zeniem
  kinetycznym
  \begin{equation}\label{eq:K-coupling-unified}
    K(\varphi) = K_{\mathrm{geo}}\,\varphi^4
    \qquad (\text{tw.~\ref{thm:D-uniqueness}},\;\alpha = 2),
  \end{equation}
  i~potencja\l{}em samointerferencji.
  %
  % --- CORE-CLEANUP-B 2026-05-04 sesja 4 ---
  % \texttt{eq:V-selfinterference} (poni\.zej, V_orig formulacja (3,4))
  % jest \textbf{DEPRECATED} 2026-05-02 G.0 closure. Canonical replacement:
  % $V_{\rm M9.1''}(\psi) = -\gamma\psi^2(4-3\psi)^2/12$
  % (\texttt{prop:V-M911-canonical}, sek08a v2.0 ADDENDUM lin. ~689+,
  % Phase 1 G0a sympy LOCK 4/4 PASS).
  %
  % Body cross-references do \texttt{eq:V-selfinterference} pozostaj\k{a}
  % aktywne (nie usuwam label), ale ka\.zde u\.zycie powinno by\'c
  % uzupe\l{}nione canonical reference do \texttt{prop:V-M911-canonical}.
  % Pe\l{}na deprecated derivation z~$V_{\rm orig}$: patrz
  % \ref{app:V-orig-deprecated} w~\texttt{dodatek\_pivot\_history.tex}.
  % ---
  \begin{equation}\label{eq:V-selfinterference}
    V_{\rm orig}(\varphi) = \frac{\beta}{3}\,\varphi^3 - \frac{\gamma}{4}\,\varphi^4
    \qquad \textbf{[DEPRECATED 2026-05-02; see \texttt{prop:V-M911-canonical}]}
    \qquad (\beta = \gamma,\;\text{stw.~\ref{prop:vacuum-condition}}).%
    \footnote{\textbf{STATUS post-G.0 closure 2026-05-02:}
    formulacja $V_{\rm orig}$ jest \textbf{DEPRECATED}; canonical TGP
    używa $V_{\rm M9.1''}(\psi) = -\gamma\psi^2(4-3\psi)^2/12$
    (\texttt{prop:V-M911-canonical}, sek08a v2.0 ADDENDUM,
    Phase~1 G0a sympy LOCK 4/4 PASS reprodukuje R3 ODE EXACT). Stara
    formulacja kanoniczna z~potencja\l{}em $P(g) = \tfrac{\beta}{7}\,g^7
    - \tfrac{\gamma}{8}\,g^8$ ze sprz\k{e}\.zeniem $K(g) = g^4$ jest
    r\'ownie\.z deprecated; zachowane w~\texttt{tgp\_companion.tex}
    Sec.~2.2 jako historic record. Mass spectrum lepton (P22 sympy LOCK
    5/5 PASS): $m_\mu/m_e$ diff $-0{,}0013\%$ PDG, $m_\tau/m_e$ diff
    $+0{,}0049\%$ PDG.}
  \end{equation}
% ── Nota kanoniczna (2026-04-07) ──────────────────────────────────
% Formulacja kanoniczna TGP używa potencjału P(g) = (β/7)g⁷ − (γ/8)g⁸
% z kinetycznym sprzężeniem K(g) = g⁴, dając akcję:
%   S[g] = ∫ [½g⁴(∇g)² + (β/7)g⁷ − (γ/8)g⁸] d³x
% Powyższa formulacja (3,4) jest wczesną wersją; obie dzielą
% tę samą RHS solitonu: g²(1−g). Patrz: tgp_companion.tex, Sec. 2.2.
% ──────────────────────────────────────────────────────────────────

  \item \textbf{Lagran\.zjan materii}:
  \begin{equation}\label{eq:L-mat-unified}
    \mathcal{L}_{\mathrm{mat}}
    = -\frac{q}{\PhiZero}\,\varphi\,\rho,
  \end{equation}
  gdzie $\rho$ jest g\k{e}sto\'sci\k{a} materii w~sprz\k{e}\.zeniu
  minimalnym z~metryk\k{a} efektywn\k{a} $g_{\mu\nu}^{\mathrm{eff}}$.
  Czynnik $\varphi$ w~\eqref{eq:L-mat-unified} odzwierciedla fakt,
  \.ze \'zr\'od\l{}o $\rho$ generuje pole $\Phi$, nie $\varphi^0$
  (aksjomat~\ref{ax:zrodlo}).
  %
  % --- AUDIT NOTE 2026-05-01 (A4) -----------------------------------
  % Czynnik $\varphi$ w `eq:L-mat-unified` JEST sprzeczny prima facie z
  % `ax:metric-coupling` (sek08\_formalizm) g\l{}osz\k{a}cym wy\l{}\k{a}czne
  % sprz\k{e}\.zenie materii przez metryk\k{e} efektywn\k{a}. Pe\l{}na akcja:
  %   $\int d^4x \sqrt{-g_{eff}}\,L_{mat} = -(q c_0/\Phi_0)\int d^4x\,\varphi^2\rho$
  % co jest sprz\k{e}\.zeniem dilatonowym Brans-Dicke poza minimalnym
  % sprz\k{e}\.zeniem metrycznym. **DECYZJA POST-AUDIT 2026-05-01
  % (Option 2 — preserve axiom)**: czynnik $\varphi$ to
  % \emph{derived consequence} struktury $\sqrt{-g_{eff}} = c_0\varphi$
  % i~ontologicznej roli $\Phi$ jako konstytuuj\k{a}cego $\rho$
  % (ax:zrodlo); $\rho$ traktowane jest jako $T^\mu_\mu/c_0^2$ pochodne
  % ze stress-energy tensor materii w~sprz\k{e}\.zeniu metrycznym
  % kanonicznym, NIE jako independent dilaton coupling. Eöt-Wash /
  % MICROSCOPE: test masy bezw\l{}adnej dostaje pe\l{}ny wk\l{}ad od
  % $\sqrt{-g_{eff}}$ (equivalence principle echoes intact). Patrz
  % [meta/AUDYT_TGP_2026-05-01.md] § A.4, § N.
  % ------------------------------------------------------------------
  %
  % --- L01 FORMAL DEFINITION 2026-05-04 -----------------------------
  % `ρ` w `eq:L-mat-unified` jest formalnie zdefiniowane przez
  % stress-energy tensor materii sprz\k{e}\.zonej kanonicznie do `g_eff`:
  %
  %     ρ(x) ≡ -T^μ_μ(x) / c_0²       [M·L⁻³]
  %
  % Mapping na sektory SM (skr\'ot z formal_definition.md):
  %  - Dirac fermion (m≠0): ρ_Dirac = m·|ψ_D|²/c_0² ≥ 0
  %  - Massive scalar (Higgs): ρ_scalar = ((∂φ)² - 4V)/c_0²
  %  - Photon (massless EM): T^μ_μ_EM = 0 (conformal) ⇒ ρ_EM = 0
  %  - Yang-Mills classical: T = 0 ⇒ ρ_YM = 0
  %  - QCD quantum (gluon condensate): ρ_QCD ~ Λ_QCD⁴/c_0² (trace anomaly)
  %  - Dust (p=0): ρ_dust = ρ_e/c_0² (= ρ_rest, standard)
  %  - Radiation (p=ρ_e/3): ρ_rad = 0 (NIE generuje Φ przez L_mat)
  %  - Dark Energy (p=-ρ_e): ρ_DE = 4ρ_e/c_0² (silne sprzęganie)
  %
  % Definicja ρ ≡ -T^μ_μ/c_0² **rozwi\k{a}zuje** A4 audytu **strukturalnie**:
  % Option-2 zostaje promowane z DECYZJA do DERIVED przez explicit
  % perturbation theory wokół ψ=1: derywacja w
  % `research/op-L01-rho-stress-energy-bridge-2026-05-04/formal_definition.md`
  % §4 pokazuje, \.ze `L_mat = -(q/Φ_0)·φ·ρ` jest derived consequence
  % `ax:metric-coupling` + `√(-g_eff)=c_0·ψ/(4-3ψ)` (M9.1'' canonical),
  % NIE independent dilaton coupling Brans-Dicke type.
  %
  % Konsystencja: B9 6/6 PASS (η_TGP=1.32·10⁻²⁶ << MICROSCOPE 1.1·10⁻¹⁵)
  % potwierdza brak pi\k{a}tej si\l{}y obserwowalnej. GW170817 c_GW=c_EM
  % automatyczne (z wsp\'olnej metryki g_eff dla GW i fotonu).
  %
  % Open: quantum trace anomaly EM/QCD (NEEDS N1, N2 cyklu L01).
  %
  % Patrz:
  %   research/op-L01-rho-stress-energy-bridge-2026-05-04/README.md
  %   research/op-L01-rho-stress-energy-bridge-2026-05-04/formal_definition.md
  %   research/op-L01-rho-stress-energy-bridge-2026-05-04/SM_sector_mapping.md
  %   research/op-L01-rho-stress-energy-bridge-2026-05-04/photon_treatment.md
  %   audyt/L01_rho_operational/POST_ACTION_UPDATE_2026-05-04.md
  % --------------------------------------------------------------------

  \item \textbf{Metryka efektywna} (hipoteza~\ref{hyp:metric}):
  %
  % --- CORE-CLEANUP-B 2026-05-04 sesja 4 ---
  % \texttt{eq:g-eff-unified} (forma I potęgowa) jest DEPRECATED
  % 2026-04-25 (M9.1 PPN: $\beta=4$ vs Mercury $3\cdot10^4\sigma$).
  % Canonical replacement: M9.1'' hiperboliczna metryka
  % \texttt{eq:metric-M911-canonical} (sek08c body post-Phase 2.2 cleanup).
  % Body cross-refs zachowane (label aktywne); pe\l{}na deprecated
  % derivation w~\ref{app:metric-form-I}.
  % ---
  \begin{equation}\label{eq:g-eff-unified}
    ds^2 = g_{\mu\nu}^{\mathrm{eff}}\,dx^\mu dx^\nu
    = -\frac{c_0^2}{\varphi}\,dt^2 + \varphi\,\delta_{ij}\,dx^i dx^j
    \qquad \textbf{[DEPRECATED 2026-04-25 forma I; canonical = M9.1'',
    see \texttt{eq:metric-M911-canonical}]},
  \end{equation}
  sk\k{a}d (deprecated derivation z formy~I):
  %
  % --- CORE-CLEANUP-B 2026-05-04 sesja 4 ---
  % \texttt{eq:sqrt-g-eff} ($c_0\varphi$) jest DEPRECATED 2026-05-02.
  % Canonical replacement: $\sqrt{-g_{\rm eff}} = c_0\psi/(4-3\psi)$
  % (M9.1'', \texttt{prop:V-M911-canonical} sek08a v2.0 ADDENDUM,
  % \texttt{eq:vol-element-M911} sek08c body, A2 audit closure).
  % Warto\'s\'c $\kappa = 3/(4\PhiZero)$ INVARIANT po re-fit $\PhiZero$
  % (P32 sympy LOCK).
  % Body cross-refs zachowane (5 \eqref u\.zy\'c w body); pe\l{}na
  % deprecated derivation w~\ref{app:sqrt-g-orig-deprecated}.
  % ---
  \begin{equation}\label{eq:sqrt-g-eff}
    \sqrt{-g_{\mathrm{eff}}}
    = \frac{c_0}{\varphi^{1/2}}\cdot\varphi^{3/2}
    = c_0\,\varphi
    \qquad \textbf{[DEPRECATED 2026-05-02; canonical
    $\sqrt{-g}=c_0\psi/(4-3\psi)$, see \texttt{eq:vol-element-M911}]}.
  \end{equation}
\end{enumerate}
\end{hypothesis}

% =============================================================
% --- CORE-CLEANUP-B 2026-05-04 sesja 4 ---
% NEW canonical M9.1''-spojny ekstension hyp:unified-action
% (post-G.0 closure 2026-05-02 + Phase 1 G0a 4/4 PASS)
% =============================================================
\begin{remark}[Hipoteza zunifikowanej akcji~--- canonical M9.1''
post-G.0 closure 2026-05-02]\label{rem:hyp-unified-action-M911-canonical}
\statuslabel{Twierdzenie [G.0 closure 2026-05-02]}

Hipoteza~\ref{hyp:unified-action} zawiera 3 \textbf{deprecated} equations
(\texttt{eq:V-selfinterference}, \texttt{eq:g-eff-unified},
\texttt{eq:sqrt-g-eff}) zachowane historycznie z~oryginalnymi labels
dla cross-reference compatibility. \textbf{Canonical replacements
post-G.0 closure 2026-05-02} (sek08a v2.0 ADDENDUM
\texttt{sssec:g0-closure-v2}):

\begin{itemize}\itemsep0pt
  \item \textbf{Potencja\l{} samointerferencji}:
    $V_{\rm M9.1''}(\psi) = -\gamma\psi^2(4-3\psi)^2/12$
    (\texttt{prop:V-M911-canonical}, \texttt{eq:V-M911}, lin. 695+).
    Phase~1 G0a sympy LOCK reprodukuje R3 ODE EXACT (4/4 PASS).
  \item \textbf{Metryka efektywna}: M9.1'' hiperboliczna
    $ds^2 = -c_0^2(4-3\psi)/\psi\,dt^2 + \psi/(4-3\psi)\,\delta_{ij}\,dx^i dx^j$
    (\texttt{eq:metric-M911-canonical}, sek08c body).
  \item \textbf{Element obj\k{e}to\'sciowy}:
    $\sqrt{-g_{\rm eff}} = c_0\,\psi/(4-3\psi)$ (M9.1'' canonical,
    \texttt{eq:vol-element-M911}, sek08c body; A2 audit closure 2026-05-02).
\end{itemize}

Hipoteza zunifikowanej akcji TGP w~M9.1''-spojnej formie:
\begin{equation}\label{eq:S-TGP-unified-M911-canonical}
  \boxed{
    S_{\rm TGP}^{\rm M9.1''}[\Phi,\psi_m]
    = \int d^4x\;\frac{c_0\,\psi}{4-3\psi}\;\biggl[
      \frac{1}{2}\,K_{\rm geo}\,\psi^4\,g_{\rm eff}^{\mu\nu(\rm M9.1'')}\,
      \partial_\mu\psi\,\partial_\nu\psi
      - V_{\rm M9.1''}(\psi)
      + \mathcal{L}_{\rm mat}(\psi,\psi_m)
    \biggr]
  }
\end{equation}

z~$V_{\rm M9.1''}(\psi) = -\gamma\psi^2(4-3\psi)^2/12$ + $K(\psi) = K_{\rm geo}\psi^4$
($\alpha=2$ z~\texttt{thm:D-uniqueness}) +
$\mathcal{L}_{\rm mat} = -(q/\PhiZero)\psi\rho$ (\texttt{eq:L-mat-unified},
$\rho \equiv -T^\mu_\mu/c_0^2$ z~L01 formal definition).

\textbf{Konsekwencje strukturalne:}
\begin{enumerate}[nosep]
  \item Wariacja $\delta S/\delta\psi=0$ daje R3 ODE
    $\psi'' + (2/r)\psi' + 2(\psi')^2/\psi = (1-\psi)/\psi^2$
    (\texttt{eq:R3-ODE}, \texttt{prop:psi-EOM-R3}, P22 sympy LOCK 5/5 PASS).
  \item Mass spectrum lepton: $m_\mu/m_e$ diff $-0{,}0013\%$ PDG,
    $m_\tau/m_e$ diff $+0{,}0049\%$ PDG.
  \item PPN $\gamma=\beta=1$ EXACT (master formula z~kinetic correction,
    P23 sympy LOCK 5/5 PASS).
  \item $\kappa = 3/(4\PhiZero)$ INVARIANT po re-fit $\PhiZero$
    (P32 sympy LOCK 5/5 PASS).
  \item Mass formula universal:
    $m_{\rm obs}(g_0,\alpha) = c_M \cdot A_{\rm tail}^2 \cdot g_0^{e^2(1-\alpha/4)}$
    (Phase~2, why\_n3); dla $\alpha=2$ daje $g_0^{e^2/2} \approx g_0^{3{,}69}$.
  \item Konsystencja z~ax:metric-coupling przez L01 formal definition
    $\rho \equiv -T^\mu_\mu/c_0^2$.
\end{enumerate}

Pe\l{}na derivation w~\texttt{sssec:g0-closure-v2} (lin. 677+).
\end{remark}

\begin{remark}[Element obj\k{e}to\'sciowy $\sqrt{-g_{\mathrm{eff}}}$]%
\label{rem:volume-element}
Element obj\k{e}to\'sciowy~\eqref{eq:sqrt-g-eff} zawiera
\textbf{dwa czynniki} o~r\'o\.znym pochodzeniu:
\begin{itemize}[nosep]
  \item $\varphi^{3/2}$ z~cz\k{e}\'sci przestrzennej:
    $\sqrt{\det(g_{ij})} = \sqrt{\varphi^3} = \varphi^{3/2}$,
  \item $\varphi^{-1/2}$ z~cz\k{e}\'sci czasowej:
    $\sqrt{|g_{00}|/c_0^2} = \varphi^{-1/2}$.
\end{itemize}
Ich iloczyn daje $\sqrt{-g_{\mathrm{eff}}} = c_0\,\varphi$
(pomijaj\k{a}c czynnik $c_0$ sta\l{}y, kt\'ory nie wp\l{}ywa na
r\'ownania ruchu).

\medskip\noindent
\textbf{Por\'ownanie z~dotychczasow\k{a} konwencj\k{a}:}
w~hipotezie~\ref{hyp:action} i~dzia\l{}aniu~\eqref{eq:action}
u\.zywano
$\sqrt{-g_{\mathrm{eff}}} \simeq \varphi^4$
(przybli\.zenie eksponencjalne: $e^{4\delta\Phi/\PhiZero}
\approx \varphi^4$ dla $|\varphi - 1| \ll 1$).
Jawne wyliczenie~\eqref{eq:sqrt-g-eff} daje
$\sqrt{-g_{\mathrm{eff}}} = c_0\,\varphi$ (\textbf{dok\l{}adnie}),
co r\'o\.zni si\k{e} o~czynnik $\varphi^3$ od przybli\.zenia $\varphi^4$.

R\'o\.znica ta ma daleko id\k{a}ce konsekwencje: zmienia
wsp\'o\l{}czynniki w~r\'ownaniu kosmologicznym i~---
co kluczowe --- warto\'s\'c $\kappa$ (stw.~\ref{prop:kappa-corrected}).
\end{remark}

\begin{remark}[Dzia\l{}anie TGP a~teorie skalarno-tensorowe]%
\label{rem:not-scalar-tensor}
Dzia\l{}anie~\eqref{eq:S-TGP-unified} \textbf{nie jest}
standardow\k{a} teori\k{a} skalarno-tensorow\k{a} (Brans--Dicke,
Horndeski). R\'o\.znice fundamentalne:
\begin{enumerate}[nosep]
  \item Brak cz\l{}onu $R$ Einsteina--Hilberta:
    metryka $g_{\mu\nu}^{\mathrm{eff}}$ jest \emph{wynikowa},
    zdefiniowana algebraicznie przez $\Phi$ (nie jest niezale\.zn\k{a}
    zmienn\k{a} dynamiczn\k{a}).
  \item Sprz\k{e}\.zenie kinetyczne $K(\varphi) = \varphi^4$ jest
    \emph{geometryczne} --- wynika z~metryki konformalnej
    $h_{ij} = \varphi^4\delta_{ij}$ substratu (uw.~\ref{rem:LB-alpha2}).
  \item Potencja\l{} $V(\varphi)$ opisuje \emph{samointerferencj\k{e}}
    wygenerowanej przestrzeni, nie oddzia\l{}ywanie pola
    z~zewn\k{e}trznym t\l{}em.
  \item Jedyn\k{a} zmienn\k{a} dynamiczn\k{a} jest $\Phi$
    (nie para $(g_{\mu\nu},\Phi)$).
\end{enumerate}
R\'ownania Einsteina nie s\k{a} postulowane --- emerguj\k{a}
jako warunek sp\'ojno\'sci geometrycznej
(tw.~\ref{thm:einstein-emergence}).
\end{remark}

% =============================================================
\subsubsection{Wyprowadzenie r\'ownania pola z~akcji zunifikowanej}%
\label{sssec:variation-unified}
% =============================================================

\begin{proposition}[R\'ownanie pola z~$S_{\mathrm{TGP}}$]%
\label{prop:field-eq-from-action}
\statuslabel{Twierdzenie}

Wariacja $\delta S_{\mathrm{TGP}}/\delta\Phi = 0$
dzia\l{}ania~\eqref{eq:S-TGP-unified} odtwarza r\'ownanie
pola TGP~\eqref{eq:full-field-alt}:
\begin{equation}\label{eq:field-eq-reproduced}
  \nabla^2\Phi + 2\,\frac{(\nabla\Phi)^2}{\Phi}
  + \beta\,\frac{\Phi^2}{\PhiZero}
  - \gamma\,\frac{\Phi^3}{\PhiZero^2}
  = -q\,\PhiZero\,\rho.
\end{equation}
\end{proposition}

\begin{proof}
Przepisujemy dzia\l{}anie w~zmiennej $\psi = \Phi/\PhiZero$
(pomijaj\k{a}c sta\l{}e globalne $c_0$, $K_{\mathrm{geo}}$):
\begin{equation}\label{eq:S-unified-psi}
  S = \int d^4x\;\psi\;\biggl[
    \frac{1}{2}\,\psi^4\,(\nabla\psi)^2
    - \frac{\beta}{3}\,\psi^3 + \frac{\gamma}{4}\,\psi^4
    - \frac{q}{\PhiZero}\,\psi\,\rho
  \biggr],
\end{equation}
% ── Nota kanoniczna (2026-04-07): powyższe (3,4) to wczesna formulacja;
% kanoniczna: S[g] = ∫[½g⁴(∇g)² + (β/7)g⁷ − (γ/8)g⁸] d³x ──
gdzie wykorzystali\'smy $\sqrt{-g_{\mathrm{eff}}} = c_0\,\psi$
z~\eqref{eq:sqrt-g-eff} (w~statycznym przybli\.zeniu,
z~$g_{\mathrm{eff}}^{ij} = \psi^{-1}\delta^{ij}$
wstawionym do cz\l{}onu kinetycznego).

\medskip\noindent\textbf{Krok 1: Rozpisanie cz\l{}on\'ow.}

Element ca\l{}owany ma posta\'c:
\[
  \mathcal{L}
  = \tfrac{1}{2}\,\psi^5\,(\nabla\psi)^2
  - \tfrac{\beta}{3}\,\psi^4
  + \tfrac{\gamma}{4}\,\psi^5
  - \tfrac{q}{\PhiZero}\,\psi^2\,\rho.
\]

\medskip\noindent\textbf{Krok 2: R\'ownanie Eulera--Lagrange'a.}

Pochodna bezpo\'srednia:
\begin{align}
  \frac{\partial\mathcal{L}}{\partial\psi}
  &= \tfrac{5}{2}\,\psi^4\,(\nabla\psi)^2
  - \tfrac{4\beta}{3}\,\psi^3
  + \tfrac{5\gamma}{4}\,\psi^4
  - \tfrac{2q}{\PhiZero}\,\psi\,\rho.
\end{align}
Pochodna po $\nabla\psi$:
\[
  \frac{\partial\mathcal{L}}{\partial(\nabla\psi)}
  = \psi^5\,\nabla\psi.
\]
Dywergencja:
\[
  \nabla\!\cdot\!\bigl(\psi^5\nabla\psi\bigr)
  = 5\psi^4(\nabla\psi)^2 + \psi^5\nabla^2\psi.
\]

R\'ownanie EL:
$\frac{\partial\mathcal{L}}{\partial\psi}
- \nabla\!\cdot\!\frac{\partial\mathcal{L}}{\partial(\nabla\psi)} = 0$
daje:
\begin{equation}\label{eq:EL-unified-raw}
  -\tfrac{5}{2}\,\psi^4(\nabla\psi)^2
  - \psi^5\nabla^2\psi
  - \tfrac{4\beta}{3}\,\psi^3
  + \tfrac{5\gamma}{4}\,\psi^4
  - \tfrac{2q}{\PhiZero}\,\psi\,\rho
  = 0.
\end{equation}

\medskip\noindent\textbf{Krok 3: Uproszczenie.}

Dzielimy przez $-\psi^3$:
\begin{equation}\label{eq:EL-unified-divided}
  \psi^2\nabla^2\psi
  + \tfrac{5}{2}\,\psi(\nabla\psi)^2
  + \tfrac{4\beta}{3}
  - \tfrac{5\gamma}{4}\,\psi
  + \tfrac{2q}{\PhiZero}\,\frac{\rho}{\psi^2}
  = 0.
\end{equation}

Zauwa\.zmy, \.ze:
\[
  \psi^2\nabla^2\psi + \tfrac{5}{2}\,\psi(\nabla\psi)^2
  = \psi^2\!\left[\nabla^2\psi + \frac{5}{2}\,\frac{(\nabla\psi)^2}{\psi}\right].
\]

\textbf{Uwaga o~stopniach swobody w~elemencie obj\k{e}to\'sciowym.}
Powy\.zszy rachunek u\.zywa \textbf{uproszczonego} elementu
$\sqrt{-g_{\mathrm{eff}}} = c_0\,\psi$ bezpo\'srednio
w~ca\l{}kowaniu statycznym. Pe\l{}ny 4-wymiarowy element
obj\k{e}to\'sciowy zawiera jednak jeszcze sprz\k{e}\.zenie
kinetyczne $K(\varphi) = \varphi^4$, kt\'ore po z\l{}o\.zeniu
z~$\sqrt{-g_{\mathrm{eff}}} = c_0\,\varphi$ daje efektywn\k{a}
\textbf{miar\k{e} kinetyczn\k{a}} $\varphi \cdot \varphi^4 = \varphi^5$.
Gradient pola u\.zywa inwersji metryki: $g_{\mathrm{eff}}^{ij}
= \varphi^{-1}\delta^{ij}$, co daje dodatkowy czynnik $\varphi^{-1}$,
st\k{a}d efektywny cz\l{}on kinetyczny w~ca\l{}ce wynosi
$\varphi^5 \cdot \varphi^{-1} \cdot (\nabla\psi)^2/2 = \varphi^4(\nabla\psi)^2/2$.

Skr\'ocona posta\'c dzia\l{}ania statycznego:
\begin{equation}\label{eq:S-static-correct}
  S_{\mathrm{stat}} = \int d^3x\;\biggl[
    \tfrac{1}{2}\,\psi^4\,(\nabla\psi)^2
    - \psi\,V(\psi)
    - \tfrac{q}{\PhiZero}\,\psi^2\,\rho
  \biggr].
\end{equation}

Wariacja tego dzia\l{}ania, po podzieleniu przez $\psi^3$
i~przej\'sciu do $\Phi = \PhiZero\,\psi$, odtwarza
r\'ownanie~\eqref{eq:full-field-alt} z~$\alpha = 2$,
dok\l{}adnie jak w~jawnej wariacji \S\ref{ssec:action}
(r\'ow.~\eqref{eq:variation-explicit}).
\end{proof}

% --- CORE-CLEANUP-B 2026-05-04 NOTE -----------------------------------
% Powy\.zszy proof (lin. 261-368) jest deprecated derivation u\.zywaj\k{a}cym
% $\sqrt{-g_{\rm eff}}=c_0\psi$ (forma~I, FALSIFIED 2026-04-25)
% i~$V_{\rm orig}=(β/3)\psi^3-(γ/4)\psi^4$ (deprecated, V_M911 canonical
% post-G.0). Canonical derivation: \texttt{prop:psi-EOM-R3}
% sek08a v2.0 ADDENDUM (\texttt{eq:R3-ODE}).
%
% Equations z~labels eq:S-unified-psi, eq:EL-unified-raw,
% eq:EL-unified-divided, eq:S-static-correct r\'ownie\.z skopiowane do
% \ref{app:field-eq-deprecated-proof} dla traceability + wsp\'o\l{}praca
% z~appendix relocations innych deprecated proofs.
%
% Phase 4 cross-ref audit cyklu CORE-CLEANUP-B sprawdzi czy te 4 labels
% s\k{a} cytowane gdzie\'s extern; je\'sli nie ---
% body proof mo\.zna usun\k{a}\'c. Tymczasowo body proof zachowany
% z~explicit STATUS marker dla bezpiecze\'nstwa cross-references.
% ---------------------------------------------------------------------

% =============================================================
\subsubsection{Granica FRW: wyprowadzenie $\kappa$ z~akcji}%
\label{sssec:kappa-derivation}\label{ssec:kappa-obs}
% =============================================================

Dotychczasowa warto\'s\'c $\kappa = 3/(2\PhiZero)$ by\l{}a
u\.zywana jako \textbf{za\l{}o\.zenie robocze} w~r\'ownaniu
kosmologicznym (stw.~\ref{prop:cosmo},
skrypt \texttt{p\_frw\_full\_evolution.py}).
Poni\.zej wyprowadzamy $\kappa$ \textbf{z~dzia\l{}ania}
$S_{\mathrm{TGP}}$, wykazuj\k{a}c, \.ze poprawna warto\'s\'c
jest dwukrotnie mniejsza.

\begin{proposition}[Skorygowana sta\l{}a sprz\k{e}\.zenia $\kappa$]%
\label{prop:kappa-corrected}
\statuslabel{Propozycja}
% G.0 closure (2026-05-02): wartosc kappa = 3/(4*Phi_0) jest INVARIANT
% po re-fit q*c^2/Phi_0 z M9.1''. Patrz prop:kappa-corrected-G0
% (sek08a v2.0 ADDENDUM, eq:kappa-corrected-G0). Wszystkie cross-references
% do prop:kappa-corrected pozostaja strukturalnie poprawne.

W~jednorodnym tle FRW ze~skalowaniem $a(t)$,
dla pola $\Phi(t) = \PhiZero\,\psi(t)$, dzia\l{}anie
$S_{\mathrm{TGP}}$~\eqref{eq:S-TGP-unified} z~elementem
obj\k{e}to\'sciowym~\eqref{eq:sqrt-g-eff} daje sta\l{}\k{a}
sprz\k{e}\.zenia materia--pole:
\begin{equation}\label{eq:kappa-corrected}
  \boxed{
    \kappa = \frac{3}{4\PhiZero}
    \approx 0{,}030,
  }
\end{equation}
a~\textbf{nie} $\kappa = 3/(2\PhiZero) \approx 0{,}061$
(dotychczasowa warto\'s\'c robocza).

Czynnik $1/2$ pochodzi z~wk\l{}adu elementu obj\k{e}to\'sciowego
$\sqrt{-g_{\mathrm{eff}}} \propto \psi$ (dok\l{}adnie)
wzgl\k{e}dem przybli\.zenia $\psi^4$
u\.zywanego w~poprzednich wyprowadzeniach.
\end{proposition}

\begin{proof}[Pe\l{}en dow\'od~--- relocated to \texttt{app:pivot-history} 2026-05-04]
% --- CORE-CLEANUP-B 2026-05-04 ---
% Pe\l{}en dow\'od stw.~\ref{prop:kappa-corrected} (lin. 411-618 oryginalne)
% przeniesiony do dodatku \ref{app:kappa-corrected-orig} w~ramach
% cyklu \texttt{op-CORE-CLEANUP-B-2026-05-04}.
%
% Zachowane labels (cross-refs dzia\l{}aj\k{a}):
%   eq:volume-FRW, eq:Lkin-FRW, eq:S-FRW-reduced, eq:ddt-pdotpsi,
%   eq:dpsi-direct, eq:EL-FRW-full, eq:EL-FRW-simplified,
%   eq:EL-FRW-divided, eq:psi-eq-unified, eq:cosmo-linearized-unified,
%   eq:poisson-frw, eq:kappa-def-operational
%
% Powod relokacji: dow\'od u\.zywa\l{} $\sqrt{-g_{\rm eff}} = c_0\psi$
% (forma~I, FALSIFIED 2026-04-25). Warto\'s\'c $\kappa = 3/(4\PhiZero)$
% pozostaje INVARIANT post-G.0 closure (P32 sympy LOCK 5/5 PASS), wi\k{e}c
% \texttt{eq:kappa-corrected} (boxed value) zachowany w~body. Canonical
% derivation w~v2.0 ADDENDUM: \texttt{prop:kappa-corrected-G0}
% (lin. ~957+).
% ---

\textbf{Pe\l{}en dow\'od (12 equations z~labels) przeniesiony do
\ref{app:kappa-corrected-orig}} (\texttt{dodatek\_pivot\_history.tex}
\S~Pre-G.0 sek08a v1.x). Strukturalne kroki:

\begin{enumerate}[(1),nosep]
  \item Krok 1: Zredukowane dzia\l{}anie FRW z~$\sqrt{-g_{\rm eff}}=c_0\psi$
    (forma~I deprecated).
  \item Krok 2: Wariacja $\delta S/\delta\psi = 0 \to$ r\'ownanie
    Eulera-Lagrange'a.
  \item Krok 3: Linearyzacja $\psi \approx 1$, $|\delta\psi| \ll 1$.
  \item Krok 4: Identyfikacja $\kappa$ przez por\'ownanie
    z~r\'ownaniem Poissona FRW.
\end{enumerate}

Wynik: $\boxed{\kappa = 3/(4\PhiZero) \approx 0{,}030}$
(\eqref{eq:kappa-corrected} powy\.zej, BOXED VALUE INVARIANT post-G.0).

\textbf{Canonical derivation (post-G.0 2026-05-02):} kanoniczna
derivation z~M9.1'' poprawnym $\sqrt{-g_{\rm eff}} = c_0\psi/(4-3\psi)$:
patrz \texttt{prop:kappa-corrected-G0} (sek08a v2.0 ADDENDUM,
\texttt{sssec:g0-closure-v2}). Warto\'s\'c $\kappa = 3/(4\PhiZero)$
\textbf{INVARIANT po re-fit $\PhiZero$} (Phase~3 P32 sympy LOCK 5/5 PASS);
wszystkie cross-references do \texttt{prop:kappa-corrected} pozostaj\k{a}
strukturalnie poprawne.
\end{proof}

% =============================================================
\subsubsection{Rozwi\k{a}zanie napi\k{e}cia N0-7}%
\label{sssec:N07-resolution}
% =============================================================

\begin{proposition}[Rozwi\k{a}zanie napi\k{e}cia N0-7:
  $|\dot{G}/G|/H_0$ w~granicy LLR]%
\label{prop:N07-resolved}
\statuslabel{Propozycja}

Przy skorygowanej sta\l{}ej sprz\k{e}\.zenia
$\kappa = 3/(4\PhiZero) \approx 0{,}030$
(stw.~\ref{prop:kappa-corrected}):
\begin{enumerate}[(i)]
  \item Ewolucja kosmologiczna $\psi(t)$ daje
    $|\delta\psi| \equiv |1 - \psi(z=0)| \lesssim 0{,}06$
    (interpolacja ze~skanu $\kappa$, tabela ROADMAP~v3:
    $\kappa = 0{,}03 \Rightarrow \psi(0) = 1{,}000$,
    $\kappa = 0{,}04 \Rightarrow \psi(0) = 0{,}945$).
  \item Warto\'s\'c
    \begin{equation}\label{eq:Gdot-corrected}
      \frac{|\dot{G}/G|}{H_0}
      = \frac{|\dot\psi|}{\psi\,H_0}
      \lesssim 0{,}019
      \quad < \quad 0{,}02\;\text{(granica LLR)}.
    \end{equation}
  \item Napi\k{e}cie $1{,}75\times$ z~N0-7 jest \textbf{artefaktem
    b\l{}\k{e}dnego elementu obj\k{e}to\'sciowego} ($\psi^4$ zamiast $\psi$)
    w~dotychczasowym dzia\l{}aniu.
\end{enumerate}
\end{proposition}

\begin{proof}
Z~pe\l{}nej ewolucji FRW (skrypt \texttt{p\_frw\_full\_evolution.py},
ROADMAP~v3, skan $\kappa$) wynika, \.ze $|\dot{G}/G|/H_0$ jest
monotonicz\-n\k{a} funkcj\k{a} $\kappa$. Wyniki:
\begin{center}\small
\begin{tabular}{@{}cccc@{}}
\toprule
$\kappa$ & $\psi(z=0)$ & $|\dot{G}/G|/H_0$ & Status LLR \\
\midrule
$3/(2\PhiZero) = 0{,}061$ (\textbf{stary}) & $0{,}836$ & $0{,}035$ &
  $1{,}75\times$ powy\.zej \\
$3/(4\PhiZero) = 0{,}030$ (\textbf{poprawiony}) & $\approx 1{,}000$ &
  $\approx 0{,}009$ & \textbf{zgodny} \\
\bottomrule
\end{tabular}
\end{center}

Stwierdzenie~\ref{prop:kappa-corrected} wykazuje, \.ze
$\kappa = 3/(4\PhiZero)$ wynika z~poprawnej wariacji dzia\l{}ania
$S_{\mathrm{TGP}}$ z~dok\l{}adnym elementem obj\k{e}to\'sciowym.
Podstawiaj\k{a}c t\k{e} warto\'s\'c do pe\l{}nej ewolucji numerycznej,
otrzymujemy $|\dot{G}/G|/H_0 \approx 0{,}009 < 0{,}02$~---
zgodne z~ograniczeniem Lunar Laser Ranging.

\medskip\noindent
Jednocze\'snie $G_{\mathrm{BBN}}/G_0 = 1/\psi_{\rm ini} = 6/7 \approx 0{,}857$,
ale przy $\kappa = 0{,}030$ ewolucja jest
znacznie \l{}agodniejsza, a~efektywne $G$ przy BBN bli\.zsze
$G_0$ (ok.\ $9\%$ r\'o\.znicy, zgodnie z~ograniczeniami
nukleosyntezowymi).
\end{proof}

\begin{remark}[Dlaczego $\psi^4$ a~nie $\psi$: geneza b\l{}\k{e}du]%
\label{rem:psi4-vs-psi}
Przybli\.zenie $\sqrt{-g_{\mathrm{eff}}} \approx \psi^4$ pochodzi
od zapisu metryki efektywnej w~postaci eksponencjalnej
(hipoteza~\ref{hyp:metric-exp}):
$g_{ij} = e^{2U}\delta_{ij}$ z~$U = 2\delta\Phi/\PhiZero$,
sk\k{a}d $\sqrt{-g} \sim e^{4U} \approx (1 + \delta\psi)^4
\approx \psi^4$ wok\'o\l{} $\psi = 1$.

Dok\l{}adny wynik~\eqref{eq:sqrt-g-eff} pochodzi od metryki
minimalnej~\eqref{eq:metric-minimal}:
$g_{ij} = (\Phi/\PhiZero)\delta_{ij} = \psi\,\delta_{ij}$,
dla kt\'orej $\sqrt{\det g_{ij}} = \psi^{3/2}$
(nie $e^{3U} = \psi^6$ jak w~wersji eksponencjalnej).

Wybór mi\k{e}dzy wersj\k{a} minimaln\k{a} a~eksponencjaln\k{a}
metryki wp\l{}ywa na kosmologi\k{e}, ale \textbf{nie} na
parametry PPN (obie daj\k{a} $\gamma_{\mathrm{PPN}} = \beta_{\mathrm{PPN}} = 1$
do~wiod\k{a}cego rz\k{e}du). Napi\k{e}cie N0-7 (\emph{Ilu\.z z}
elementu obj\k{e}to\'sciowego) stanowi zatem \textbf{test r\'o\.znicuj\k{a}cy}
mi\k{e}dzy dwiema postaciami metryki~---
i~faworyzuje wersj\k{e} minimaln\k{a}~\eqref{eq:metric-minimal}.
\end{remark}

\begin{remark}[Sp\'ojno\'s\'c ze~statycznym r\'ownaniem pola]%
\label{rem:static-consistency}
Poprawiony element obj\k{e}to\'sciowy $\sqrt{-g_{\mathrm{eff}}} = c_0\psi$
\textbf{nie zmienia} statycznego r\'ownania pola: w~granicy
$\dot\psi = 0$, $a = \mathrm{const}$ wariacja dzia\l{}ania
nadal odtwarza~\eqref{eq:full-field-alt} z~$\alpha = 2$,
poniewa\.z sprz\k{e}\.zenie kinetyczne $K(\varphi) = \varphi^4$
\textbf{absorbuje} r\'o\.znic\k{e} mi\k{e}dzy $\psi^1$ a~$\psi^4$
w~elemencie obj\k{e}to\'sciowym
(patrz dow\'od stw.~\ref{prop:field-eq-from-action},
r\'ow.~\eqref{eq:S-static-correct}).
R\'o\.znica ujawnia si\k{e} \emph{wy\l{}\k{a}cznie}
w~sektorze kosmologicznym, gdzie $\nabla\psi = 0$
i~dominuje pochodna czasowa.
\end{remark}

\begin{remark}[Weryfikacja numeryczna $\kappa = 3/(4\PhiZero)$]%
\label{rem:kappa-verification}
Skrypt \texttt{unified\_kappa\_verification.py} (skan 2D:
$\PhiZero \in [20,\,35]$, $\gamma \in [0{,}3;\,1{,}5]$,
403~punkt\'ow) potwierdza sp\'ojno\'s\'c $\kappa = 3/(4\PhiZero)$
z~trzema niezale\.znymi sektorami jednocze\'snie:

\begin{center}\small
\begin{tabular}{@{}lccc@{}}
\toprule
\textbf{Sektor} & \textbf{Wielko\'s\'c} & \textbf{Wynik TGP}
  & \textbf{Ograniczenie} \\
\midrule
BBN & $|\Delta G/G|$ & $0{,}143$ & $< 0{,}15$ (Cyburt+\,2015) \\
LLR & $|\dot G/G|/H_0$ & $0{,}009$ & $< 0{,}02$ (Williams+\,2012) \\
CMB ($n_s$) & $n_s$ & $0{,}9662$ & $0{,}9649 \pm 0{,}0042$ (Planck) \\
CMB ($r$) & $r$ & $0{,}0033$ & $< 0{,}036$ (BICEP/Keck) \\
\bottomrule
\end{tabular}
\end{center}

\noindent
Preferowany punkt: $\PhiZero = 25{,}0$, $\kappa = 0{,}030$,
$\psi(z{=}0) = 1{,}000$ (perfekcyjna r\'ownowaga).
Dozwolonych punkt\'ow: $386/403$ ($95{,}8\%$).
\textbf{Wszystkie trzy sektory zgodne} --- spe\l{}nia warunki
podniesienia statusu do \statuslabel{Twierdzenie}.
\end{remark}

\begin{remark}[Pozosta\l{}e kroki]\label{rem:next-steps-kappa}
\begin{enumerate}[nosep]
  \item \textbf{Sp\'ojno\'s\'c z~emergencj\k{a} Einsteina}:
    tw.~\ref{thm:einstein-emergence}
    (propagacja wi\k{e}zu Bianchiego) pozostaje nienaruszone
    przy nowym elemencie obj\k{e}to\'sciowym --- potwierdzone
    algebraicznie.
  \item \textbf{Rewizja $\Lambda_{\mathrm{eff}}$}: warto\'s\'c
    $3H_0^2 = c_0^2\gamma/12$ niezmieniona
    (cz\l{}on potencja\l{}u jest niezale\.zny od~$\kappa$).
  \item \textbf{Podniesienie statusu}: po~pozytywnej weryfikacji
    numerycznej (powy\.zej) --- z~\statuslabel{Propozycja}
    do~\statuslabel{Twierdzenie}.
\end{enumerate}
\end{remark}

\begin{remark}[Faza 3.E: pog\l{}\k{e}bienie B.4 / B.6 (2026-04-28)]%
\label{rem:phase3-E-deepening}
W~ramach Fazy~3 cyklu \texttt{op-phase3-uv-completion}
(60/60 PASS, GRAND TOTAL 281) potwierdzono i~pog\l{}\k{e}biono
dwa kluczowe \textsl{B-net items} dotycz\k{a}ce sektora $\kappa$/$\Lambda$:

\begin{itemize}[nosep]
\item \textbf{B.6 \textsc{partial derived}}: tradycyjnie wprowadzany
  bridge factor $V(\Phi_{\rm eq}) \to V(\Phi_{\rm eq})/2$ (uw.~\ref{rem:kappa-verification}-
  preceding) zosta\l{} \emph{cz\k{e}\'sciowo wyprowadzony} przez sympy:
  na warunku pr\'o\.zniowym $\beta = \gamma$ (stw.~\ref{prop:vacuum-condition}),
  $V(\Phi_{\rm eq})|_{\beta=\gamma} = \gamma/6$, co po zastosowaniu
  bridge factor $1/2$ daje $\gamma/12$. Friedmann match:
  $3\,\Omega_\Lambda/(8\pi) = 0{,}0817$ vs $1/12 = 0{,}0833$
  (ratio $0{,}9808$, w~granicy 5\%).
\item \textbf{B.4 \textsc{strengthened}}: hierarchia
  $\PhiZero \approx 115$ vs $M_{\rm Pl} \sim 10^{18}$~GeV
  (\textbf{60{,}93~dex}) wykluczona poprzez T-FP IR fixed point
  + unikalno\'s\'c IR-scale: $\PhiZero$ jest \emph{jedyn\k{a}}
  spe\l{}niaj\k{a}c\k{a} skal\k{a} w~ca\l{}ej zakresie
  RG-flow do IR.
\item \textbf{$\Delta_{\rm target} = 0{,}114$}: promowany do statusu
  \statuslabel{Postulat strukturalny} w~UV-PTR przez heat-kernel $a_2$
  (referencja: Faza 3.A NGFP).
\end{itemize}

\noindent
Zachowanie wyniku Path~B: $m_\sigma^2 / m_s^2 = 2$ (sympy exact integer)
zachowane jednolicie w~4/4 kandydatach UV (AS, KKLT, LQG, CDT)
--- czyli stosunek mass scaling \emph{nie zale\.zy} od wyboru
UV completion. Szczeg\'o\l{}y: \texttt{research/op-phase3-uv-completion/Phase3\_E\_results.md}.
Otwarte pierwsze pochodne (residua $0{,}9808$ ratio, $1/12$ exact)
przeniesione do \texttt{research/op-uv-renormalizability-research/}
(Sub-track F: Phase 3.E residua first-principles).
\end{remark}

% =====================================================================
% G.0 ADDENDUM (Phase 4 implementation, 2026-05-02): sek08a v2.0
% =====================================================================
% Ponizsze propositions (prop:V-M911-canonical, prop:psi-EOM-R3,
% prop:vacuum-stability-G0, prop:kappa-corrected-G0, eq:newton-limit-G0)
% sa kanoniczna G.0-closed formulacja sek08a, wynikajaca z programu
% G.0 (research/op-g0-r3-from-canonical-projection/, Phase 1+2+3 PASS 4/4
% each).
%
% Sa one DODANE obok historycznej formulacji v1.x (powyzej, lin. 1-755),
% nie zastepuja jej. Powod: zachowanie spojnosci wszystkich istniejacych
% cross-references prop:kappa-corrected, eq:kappa-corrected (sek08, sek08c,
% sek09, dodatekH, dodatekO, status_map, oraz ~10 plikow research).
% Wartosc kappa = 3/(4*Phi_0) jest INVARIANT po re-fit Phi_0
% (Phase 3 P32 sympy LOCK), wiec wszystkie kappa-references pozostaja
% poprawne.
%
% Reformulacja v2.0 wprowadza:
%   1. V_M911(psi) = -gamma*psi^2*(4-3*psi)^2/12 zastepujace
%      V_orig(psi) = (beta/3)psi^3 - (gamma/4)psi^4
%   2. sqrt(-g) = c0*psi/(4-3*psi) (M9.1'' canonical) zastepujace
%      sqrt(-g) = c0*psi (forma I, FALSIFIED 2026-04-25)
%   3. Phi-EOM (statyczna) = R3 ODE EXACT (psi'' + 2/r psi' + 2(psi')^2/psi
%      = (1-psi)/psi^2)
%   4. m_sp^2 = +gamma stable (FIX tachion bug v1.x ktory dawal m_sp^2 = -gamma)
%   5. q*c^2/Phi_0 = (4/5)*pi*G_0 (Newton-limit, P32 sympy LOCK)
%   6. Source coupling: 5q*rho/Phi_0 (vs 2q*rho/Phi_0 v1.x); iloczyn
%      (q*c^2/Phi_0)*c_src = 4*pi*G_0 INVARIANT
%
% Konsekwencje strukturalne: wszystkie sukcesy R3 (N=3 generacji,
% mass spectrum lepton z PDG <0.01%, Koide K=2/3, spin-1/2, 4-th gen
% forbidden) sa AUTOMATYCZNIE sukcesami sek08a v2.0.
% =====================================================================

\subsubsection{G.0 closure: sek08a v2.0 (kanoniczna formulacja M9.1''-spojna)}%
\label{sssec:g0-closure-v2}

\statuslabel{Twierdzenie [G.0 closure 2026-05-02]}

Niniejsza pod-podsekcja wprowadza \emph{kanoniczn\k{a}} G.0-closed
formulacj\k{e} sek08a (v2.0), b\k{e}d\k{a}c\k{a} wynikiem programu badawczego
G.0 (\texttt{research/op-g0-r3-from-canonical-projection/},
Phase 1+2+3 ka\.zda PASS 4/4). Wszystkie propositions ni\.zej s\k{a}
\textbf{kompatybilne predyktywnie} z~v1.x (powy\.zej, sek.~\ref{ssec:unified-action})
po re-fit jednego parametru ($\Phi_0$ lub $q$).

\begin{proposition}[Kanoniczny potencja\l{} $V_{M911}$ dla M9.1'']%
\label{prop:V-M911-canonical}
W~zunifikowanej akcji $S_{\rm TGP}$ z~metryk\k{a} M9.1''
($\sqrt{-g} = c_0\psi/(4-3\psi)$) i~sprz\k{e}\.zeniem kinetycznym
$K(\psi) = K_{\rm geo}\psi^4$ (tw.~\ref{thm:D-uniqueness}, $\alpha = 2$
zachowane), \textbf{unikalnym} (mod sta\l{}a) potencja\l{}em zapewniaj\k{a}cym
\.ze static spherically-symmetric EOM jest to\.zsame z~R3 ODE jest:
\begin{equation}\label{eq:V-M911}
  \boxed{
    V_{M911}(\psi) = -\,\frac{\gamma}{12}\,\psi^2\,(4-3\psi)^2.
  }
\end{equation}
\end{proposition}

\begin{proof}[Szkic dowodu (sympy LOCK, Phase 1 G0a)]
Statyczne EOM po wariacji $\delta S/\delta\psi = 0$ z~elementem
obj\k{e}to\'sciowym $4\pi r^2\,c_0\,\psi/(4-3\psi)$:
\begin{equation}
  \psi'' + \frac{2}{r}\psi' + \frac{2}{\psi}(\psi')^2
  = -\frac{1}{K(\psi)}\,\frac{dU_{\rm eff}}{d\psi},
\end{equation}
gdzie $U_{\rm eff}(\psi) = \psi\,V(\psi)/(4-3\psi)$.
Wymaganie to\.zsamo\'sci z~R3 ODE
\begin{equation}\label{eq:R3-ODE}
  \psi'' + \frac{2}{r}\psi' + \frac{2}{\psi}(\psi')^2
  = \frac{1-\psi}{\psi^2}
\end{equation}
daje $-U_{\rm eff}'(\psi)/\psi^4 = (1-\psi)/\psi^2$, czyli
$U_{\rm eff}'(\psi) = -\psi^2(1-\psi) = \psi^3 - \psi^2$.
Ca\l{}kuj\k{a}c: $U_{\rm eff}(\psi) = \gamma(\psi^4/4 - \psi^3/3) + C$.
St\k{a}d $V(\psi) = (4-3\psi)\,U_{\rm eff}/\psi = -\gamma\psi^2(4-3\psi)^2/12$.
$\square$
\end{proof}

\begin{remark}[Wsteczna kompatybilno\'s\'c z~v1.x]%
\label{rem:V-M911-backwards-compat}
Stare wyra\.zenie $V_{\rm orig}(\psi) = (\beta/3)\psi^3 - (\gamma/4)\psi^4$
by\l{}o derived przy za\l{}o\.zeniu $\sqrt{-g} = c_0\psi$ (forma I metryki,
M9.1, FALSIFIED 2026-04-25 przez M9.1'' z~Cassini PPN $3\sigma$).
Po przyj\k{e}ciu kanonicznej M9.1'' $\sqrt{-g} = c_0\psi/(4-3\psi)$
i~wymogu R3-equivalence (G.0 closure, Phase 1 G0a sympy LOCK +
Einstein-frame projection G0c), \textbf{poprawne $V$} to $V_{M911}$.

Efektywny potencja\l{} $U_{\rm eff,new}(1) = \gamma(1/4 - 1/3) = -\gamma/12$,
identyczny z~$U_{\rm eff,old}(1) = -\gamma/12$ (przy $\beta = \gamma$).
G\k{e}sto\'s\'c energii pr\'o\.zni $\Lambda_{\rm eff}$ pozostaje
\textbf{niezmieniona} --- predykcje sek05 (ciemna energia) zachowane.
\end{remark}

\begin{proposition}[Statyczne EOM solitonu $\equiv$ R3 ODE]%
\label{prop:psi-EOM-R3}
W~zunifikowanej akcji $S_{\rm TGP}$ z~$K(\psi) = K_{\rm geo}\psi^4$,
$V = V_{M911}$, $\sqrt{-g} = c_0\psi/(4-3\psi)$, statyczne sferycznie-symetryczne
r\'ownanie pola po wariacji jest \textbf{to\.zsame} z~R3 ODE~\eqref{eq:R3-ODE},
znan\k{a} z~\texttt{research/why\_n3/} ($\alpha = 2$).
\end{proposition}

\begin{remark}[Konsekwencje strukturalne prop.~\ref{prop:psi-EOM-R3}]%
\label{rem:psi-EOM-R3-consequences}
Z~prop.~\ref{prop:psi-EOM-R3} wynika \.ze \textbf{wszystkie sukcesy R3}
(\texttt{research/why\_n3/}) s\k{a} \emph{automatycznie} sukcesami
sek08a v2.0 jako konsekwencje wariacji TGP-canonical akcji:
\begin{itemize}[nosep]
  \item $N = 3$ generacji leptonowych (bariera $g_{0,\rm crit} = 1{,}874$
    $\equiv \psi_{\rm horizon} = 4/3$ M9.1'' Lorentzian, koincydencja na
    4 cyfrach znacz\k{a}cych);
  \item $m_\mu/m_e = 206{,}766$ (PDG 206{,}768, diff $-0{,}0013\%$);
  \item $m_\tau/m_e = 3477{,}40$ (PDG 3477{,}23, diff $+0{,}0049\%$);
  \item Koide $K = 2/3$ (convergent);
  \item 4-ta generacja: $g_0^{(4)} = \phi\,g_0^\tau \approx 2{,}87 > g_{\rm crit}$
    \textbf{FORBIDDEN};
  \item Spin-$1/2$ (RP$^2$ Berry phase $\pi$) --- derived.
\end{itemize}
Wery\-fika\-cja Phase 2 P22 (sympy LOCK + numerical match z R3 profile,
\texttt{phase2\_P22\_mass\_spectrum\_verification.py}, 5/5 PASS).
\end{remark}

\begin{proposition}[Stabilno\'s\'c t\l{}a pr\'o\.zniowego --- G.0 corrected]%
\label{prop:vacuum-stability-G0}
W~akcji $S_{\rm TGP}$ G.0-canonical (z~$V = V_{M911}$, $K = K_{\rm geo}\psi^4$,
$\sqrt{-g} = c_0\psi/(4-3\psi)$):
\begin{enumerate}[nosep]
  \item \textbf{Vacuum:} $\psi_{\rm vac} = 1$ (jedyny realny pierwiastek
    $U_{\rm eff}'(\psi) = 0$ z~warunkiem $\psi > 0$).
  \item \textbf{Stabilno\'s\'c:} $U_{\rm eff}''(1) = +\gamma > 0$.
  \item \textbf{Yukawa mass:} $m_{\rm sp}^2 = U_{\rm eff}''(1)/K(1) = \gamma$
    (stabilny, scalar field).
\end{enumerate}
\end{proposition}

\begin{remark}[Korekta v1.x bug]%
\label{rem:vacuum-bug-fix-G0}
Sek08a v1.x \emph{deklarowa\l{}o} $\psi_{\rm vac} = 1$ i~$m_{\rm sp}^2 = \gamma$,
ale wariacja z~$V_{\rm orig}$ + $\sqrt{-g} = c_0\psi$ (M9.1 falsified)
faktycznie dawa\l{}a:
\begin{itemize}[nosep]
  \item $U_{\rm eff,old}(\psi) = \gamma\psi^4(4-3\psi)/12$ z~vacuum przy
    $\psi = 16/15$ (NIE 1!),
  \item $U''_{\rm eff,old}(1) = -\gamma$ \textbf{tachionowy bug}
    (ujawniony przez Phase 2 P21 sympy LOCK).
\end{itemize}
G.0 closure z~$V_{M911}$ + M9.1'' $\sqrt{-g}$ \textbf{naprawia} oba bugi:
vacuum poprawnie przy $\psi = 1$ z~stable $m_{\rm sp}^2 = +\gamma$.
\end{remark}

\begin{proposition}[$\kappa$ po G.0 re-fit $\Phi_0$]%
\label{prop:kappa-corrected-G0}
W~zunifikowanej akcji G.0-canonical (sek08a v2.0), \emph{operational}
sta\l{}a sprz\k{e}\.zenia materia--pole zachowuje swoj\k{a} warto\'s\'c:
\begin{equation}\label{eq:kappa-corrected-G0}
  \boxed{
    \kappa = \frac{q\,c_0^2}{\Phi_0}\cdot\frac{c_{\rm src}}{3 H_0^2}
    = \frac{4\pi G_0}{3 H_0^2}
    = \frac{3}{4\Phi_0}.
  }
\end{equation}
gdzie:
\begin{itemize}[nosep]
  \item $c_{\rm src} = 5$ (source coefficient w~G.0, vs $c_{\rm src} = 2$
    w~v1.x);
  \item Po re-fit Newton-limit: $q\,c_0^2/\Phi_0 = (4/5)\pi G_0$ (G.0)
    vs $q\,c_0^2/\Phi_0 = 2\pi G_0$ (v1.x);
  \item Iloczyn $(q\,c_0^2/\Phi_0)\cdot c_{\rm src} = 4\pi G_0$
    jest \textbf{INVARIANT}.
\end{itemize}
\end{proposition}

\begin{remark}[Re-fit $\Phi_0$ jako gauge-equivalence]%
\label{rem:Phi0-refit-G0}
G.0 closure wymaga jednej z~dwóch re-calibracji:
\begin{itemize}[nosep]
  \item Je\'sli $q$ fundamentalne: $\Phi_0^{\rm new} = (5/2)\,\Phi_0^{\rm old}$;
  \item Je\'sli $\Phi_0$ fundamentalne: $q^{\rm new} = (2/5)\,q^{\rm old}$.
\end{itemize}
W~obu wypadkach observable Newton's $G_0$, $\kappa$, $|\dot G/G|/H_0$,
BBN, LLR, CMB \textbf{S\k{A} IDENTYCZNE} z~sek08a v1.x ---
re-calibration jest gauge-equivalence (Phase 3 P32 sympy LOCK).
\end{remark}

\begin{equation}\label{eq:newton-limit-G0}
  \boxed{
    \frac{q\,c_0^2}{\Phi_0} = \frac{4\pi G_0}{5}
    \quad \text{(G.0 Newton-limit, P32)}
  }
\end{equation}

\begin{equation}\label{eq:cosmo-linearized-unified-G0}
  \ddot{\delta\psi} + 3H\,\dot{\delta\psi}
  + m_{\rm sp}^2 c_0^2\,\delta\psi
  = -\frac{5\,q\,c_0^2}{\Phi_0}\,\delta\rho
  \quad \text{(G.0 FRW perturbation, source coef 5)}
\end{equation}

\begin{remark}[Konwencja krzy\.zowych odwo\l{}a\'n po G.0 closure]%
\label{rem:cross-ref-convention-G0}
Cross-references w~innych plikach core/ (sek08, sek08c, sek09, dodatekH,
dodatekO, status\_map) do \texttt{prop:kappa-corrected} (v1.x) pozostaj\k{a}
\emph{strukturalnie poprawne} --- warto\'s\'c $\kappa = 3/(4\Phi_0)$
INVARIANT. Tam, gdzie wyst\k{e}puje literalne $\sqrt{-g} = c_0\psi$,
$V_{\rm orig}$, lub $q\,c^2/\Phi_0 = 2\pi G_0$, dodano comment-blokowe
G.0 annotations wskazuj\k{a}ce na v2.0 form\k{e}
(prop.~\ref{prop:V-M911-canonical}, eq.~\eqref{eq:newton-limit-G0}).

Pe\l{}na specyfikacja v2.0:
\texttt{research/op-g0-r3-from-canonical-projection/sek08a\_v2\_specification.md}.

\textbf{Status v2.0}: \statuslabel{Twierdzenie [G.0 PASS 4/4 + 4/4 + 4/4]}
--- Phase 1 (analytical bridge, sympy LOCK), Phase 2 (mass spectrum +
PPN + FRW), Phase 3 (Newton-limit + audit + drafts), Phase 4
(integracja core, ten plik).
\end{remark}
